\section{Positive but Not Completely Positive Maps}
\label{ch:positive_not_cp_supporting_results}
This section proves the positive-filter identities used in trace
normalization, the Schmidt-rank linear algebra behind the maximal-overlap
theorem, the extremal projection bounds in Ky Fan's maximum principle, the
right-tensor identities behind the Choi-compression criteria, and the
elementary closure properties of $k$-positive maps earlier in this chapter.
The final sections give the convex geometry of Schmidt-number sets and the
Choi trace pairing used for entanglement detection.

\section{Positive filters and trace normalization}

This section proves the normalization and invertibility steps used in
Lemma~\ref{lem:making_positive_maps_trace_preserving}. Here $T^*$ denotes the
adjoint for the trace pairing, $\Id$ is the identity matrix, and $A^{-1/2}$
denotes the inverse of the positive square root of a positive definite matrix.

\begin{theorem}[Conjugation filters preserve positivity]
    \label{thm:conjugation_filters_preserve_positivity}
    \lean{unitaryConjLM_isCPMap, unitaryConjLM_isPositiveMap,
        IsPositiveMap.comp_unitaryConjLM, unitaryConjLM_isTP_of_unitary,
        unitaryConjLM_isChannel_of_unitary}
    \leanok
    \uses{def:cp_map, def:positive_map, def:trace_preserving, def:channel}
    For every $X\in\MN{D}$, the conjugation filter
    $\rho\mapsto X\rho X^\dagger$ is completely positive, hence positive.
    Therefore, if $T:\MN{D}\to\MN{D}$ is positive, then
    $\rho\mapsto T(X\rho X^\dagger)$ is positive. If in addition
    $X^\dagger X=\Id$, then the conjugation filter is trace-preserving and
    hence is a quantum channel.
\end{theorem}

\begin{proof}\leanok
    Complete positivity follows from the Kraus representation
    (\ref{eq:channel_kraus_representation}) with the single Kraus operator
    $X$. Positivity follows immediately. If $X^\dagger X=\Id$, cyclicity of
    the trace gives
    $\tr(X\rho X^\dagger)=\tr(X^\dagger X\rho)=\tr(\rho)$.
\end{proof}

\begin{lemma}[Trace-normalization criterion, trace-preservation component]
    \label{lem:trace_normalization_criterion_trace}
    \lean{trace_comp_unitaryConjLM_of_conj_traceAdjointMap_one}
    \leanok
    \uses{def:trace_preserving_rectangular, def:trace_pairing_adjoint}
    Let $T:\MN{D}\to\MN{D'}$ be a linear map and let $X\in\MN{D}$ satisfy
    $X^\dagger T^*(\Id)X=\Id$. Then, for every $\rho\in\MN{D}$,
    \begin{align}
        \tr\!(T(X\rho X^\dagger))
         & =\tr(\rho).
        \notag
    \end{align}
\end{lemma}

\begin{proof}\leanok
    Expanding with the trace-pairing adjoint gives
    \begin{align}
        \tr\!(T(X\rho X^\dagger))
         & =\tr\!(X\rho X^\dagger T^*(\Id))
        =\tr\!(\rho X^\dagger T^*(\Id)X)
        =\tr(\rho).
        \notag
    \end{align}
\end{proof}

\begin{theorem}[Filtered trace-normalization criterion]
    \label{thm:filtered_trace_normalization_criterion}
    \lean{comp_unitaryConjLM_positive_tracePreserving_of_conj_traceAdjointMap_one}
    \leanok
    \uses{def:positive_map_rectangular, def:trace_preserving_rectangular,
        def:trace_pairing_adjoint, lem:trace_normalization_criterion_trace,
        thm:conjugation_filters_preserve_positivity}
    Let $T:\MN{D}\to\MN{D'}$ be a positive map between matrix algebras of
    possibly different dimensions, and let $X\in\MN{D}$ satisfy
    $X^\dagger T^*(\Id)X=\Id$, where $T^*$ is the trace-pairing adjoint.
    Then $\rho\mapsto T(X\rho X^\dagger)$ is positive and trace-preserving.

    This is the algebraic trace-normalization step in the proof of
    \cite[Chapter~3, Lemma ``Making positive maps trace preserving'']
    {Wolf2012Quantum}. The inverse-square-root choice is recorded in the
    next theorem.
\end{theorem}

\begin{proof}\leanok
    By Theorem~\ref{thm:conjugation_filters_preserve_positivity},
    $X\rho X^\dagger$ is positive whenever $\rho$ is positive; positivity
    of $T$ therefore makes the filtered map positive. For trace preservation, the
    trace-pairing adjoint identity and cyclicity of the trace give
    \begin{align}
        \tr\!(T(X\rho X^\dagger))
         & =\tr\!(T^*(\Id)X\rho X^\dagger)
        =\tr\!(X^\dagger T^*(\Id)X\rho)
        =\tr(\rho).
        \notag
    \end{align}
\end{proof}

\begin{lemma}[Normalizing conjugation by inverse square root]
    \label{lem:conj_inv_sqrt_normalizes}
    \lean{Matrix.conjTranspose_inv_sqrt_mul_self_mul_inv_sqrt_eq_one_of_posDef}
    \leanok
    \uses{}
    Let $A\in\MN{D}$ be positive definite and set $S=A^{1/2}$. Then
    \begin{align}
        (S^{-1})^\dagger A S^{-1}
         & =\Id.
        \notag
    \end{align}
\end{lemma}

\begin{proof}\leanok
    Since $A$ is positive definite, $S$ is invertible and self-adjoint, and
    $S^2=A$. Therefore
    $(S^{-1})^\dagger A S^{-1}=S^{-1}S^2S^{-1}=\Id$.
\end{proof}

\begin{lemma}[Inverse square root of a positive definite matrix is invertible]
    \label{lem:pos_def_inv_sqrt_invertible}
    \lean{Matrix.PosDef.isUnit_det_inv_sqrt}
    \leanok
    \uses{}
    If $A\in\MN{D}$ is positive definite, then $(A^{1/2})^{-1}$ is
    invertible.
\end{lemma}

\begin{proof}\leanok
    A positive definite matrix is invertible, hence so is its square root
    $A^{1/2}$, and the inverse of an invertible matrix is invertible.
\end{proof}

\begin{theorem}[Trace normalization by inverse square root]
    \label{thm:trace_normalization_inverse_square_root}
    \lean{comp_unitaryConjLM_inv_cfc_sqrt_traceAdjointMap_one}
    \leanok
    \uses{def:positive_map_rectangular, def:trace_preserving_rectangular,
        def:trace_pairing_adjoint,
        thm:filtered_trace_normalization_criterion, lem:conj_inv_sqrt_normalizes}
    Let $T:\MN{D}\to\MN{D'}$ be a positive map such that $T^*(\Id)$ is positive
    definite. Put $X=(T^*(\Id))^{-1/2}$. Then
    $\rho\mapsto T(X\rho X^\dagger)$ is positive and trace-preserving.
\end{theorem}

\begin{proof}\leanok
    Let $S=(T^*(\Id))^{1/2}$. Since $T^*(\Id)$ is positive definite, $S$ is
    invertible and self-adjoint, and $S^2=T^*(\Id)$. For $X=S^{-1}$,
    \begin{align}
        X^\dagger T^*(\Id)X
         & =S^{-1}S^2S^{-1}=\Id.
        \notag
    \end{align}
    Theorem~\ref{thm:filtered_trace_normalization_criterion} applies.
\end{proof}

Consequently, the inverse-square-root filter used in
Lemma~\ref{lem:making_positive_maps_trace_preserving} is invertible and
satisfies the required normalization identity.


\section{Schmidt-rank factorization and spectral expansions}
\label{sec:app_positive_not_cp_schmidt_rank}

This section supplies the bounded-rank factorizations and spectral identities
used in Theorems~\ref{thm:maximal_overlap_schmidt_rank},
\ref{thm:spectral_lower_bound_schmidt_rank}, and
\ref{thm:spectral_upper_bound_schmidt_rank}. The coefficient matrix of a
bipartite vector identifies Schmidt rank with ordinary matrix rank; reduced
densities and eigenbasis expansions then express the relevant overlaps and
quadratic forms.

\begin{lemma}[Monotonicity of bounded Schmidt rank]\label{lem:schmidt_rank_mono}
    \lean{Matrix.HasSchmidtRankLE.mono}
    \leanok
    \uses{def:schmidt_rank}
    If $\SR(\psi)\le k$ and $k\le l$, then $\SR(\psi)\le l$.
\end{lemma}

\begin{proof}
    \leanok
    The hypothesis gives $\SR(\psi)\le k$, and transitivity with $k\le l$
    gives $\SR(\psi)\le l$.
\end{proof}

\begin{lemma}[Rank factorization through a coordinate space]
    \label{lem:matrix_rank_factorization_coordinate}
    \lean{Matrix.exists_mul_eq_of_rank_le}
    \leanok
    If $A\in\MN{m,n}(\C)$ has rank at most $k$, then there are matrices
    $B\in\MN{m,k}(\C)$ and $C\in\MN{k,n}(\C)$ such that $BC=A$.
\end{lemma}

\begin{proof}
    \leanok
    The range of the linear map represented by $A$ has dimension at most $k$.
    Embed this range into $\C^k$, compose with the range map, and then map
    back to the ambient codomain. The corresponding matrices give the desired
    factorization.
\end{proof}

\begin{lemma}[Coefficient-matrix form of the reduced density]
    \label{lem:reduced_eig_density_factored}
    \lean{Matrix.IsHermitian.reducedEigDensity_eq}
    \leanok
    \uses{def:reduced_eig_density}
    Writing $C_i$ for the coefficient matrix of $\phi_i$, the reduced density
    operator factors as $\rho_i=C_iC_i^\dagger$.
\end{lemma}

\begin{proof}
    \leanok
    The partial trace of $\ketbra{\phi_i}{\phi_i}$ over the second factor has
    entries $\sum_b(C_i)_{a,b}\overline{(C_i)_{a',b}}$, which is the
    $(a,a')$ entry of $C_iC_i^\dagger$.
\end{proof}

\begin{lemma}[Reduced density is positive semidefinite]
    \label{lem:reduced_eig_density_posSemidef}
    \lean{Matrix.IsHermitian.reducedEigDensity_posSemidef}
    \leanok
    \uses{lem:reduced_eig_density_factored}
    The reduced density operator $\rho_i$ is positive semidefinite.
\end{lemma}

\begin{proof}
    \leanok
    By the coefficient-matrix form, $\rho_i=C_iC_i^\dagger$, which is positive
    semidefinite for any matrix $C_i$.
\end{proof}

\begin{lemma}[Eigenbasis Rayleigh expansion]
    \label{lem:eigenbasis_rayleigh}
    \lean{Matrix.IsHermitian.rayleigh}
    \leanok
    For a Hermitian operator $\tau$ with eigenvalues $\nu_i$ and normalized
    eigenvectors $\phi_i$, the expectation in a vector $\psi$ expands as
    \begin{align}
        \bra{\psi}\tau\ket{\psi}
         & =\sum_i\nu_i|\braket{\phi_i}{\psi}|^2.
        \label{eq:schwarz_eigenbasis_rayleigh}
    \end{align}
\end{lemma}

\begin{proof}
    \leanok
    Diagonalizing $\tau=U\diag(\nu_i)U^\dagger$ with $U$ the unitary of
    eigenvectors, the change of variable $y=U^\dagger\psi$ gives
    $\bra{\psi}\tau\ket{\psi}=\sum_i\nu_i|y_i|^2$, and
    $y_i=\braket{\phi_i}{\psi}$ is the $i$-th eigenvector overlap.
\end{proof}

\begin{lemma}[Eigenbasis Parseval identity]
    \label{lem:eigenbasis_parseval}
    \lean{Matrix.IsHermitian.sum_normSq_eigenvector_overlap}
    \leanok
    For a Hermitian operator $\tau$ with normalized eigenvectors $\phi_i$, the
    eigenvector overlaps recover the squared norm of $\psi$:
    \begin{align}
        \sum_i|\braket{\phi_i}{\psi}|^2
         & =\|\psi\|^2.
        \label{eq:schwarz_eigenbasis_parseval}
    \end{align}
\end{lemma}

\begin{proof}
    \leanok
    The eigenvectors form an orthonormal basis, so with $y=U^\dagger\psi$ one
    has
    $\sum_i|\braket{\phi_i}{\psi}|^2=\sum_i|y_i|^2=\|y\|^2=\|\psi\|^2$,
    where the last equality follows because $U$ is unitary.
\end{proof}

The factorization lemma gives coordinate representatives for bounded Schmidt
rank, while the reduced-density and eigenbasis identities provide the
Frobenius, Rayleigh, and Parseval formulas used in the main spectral bounds.

\section{Ky Fan's maximum principle}
\label{sec:app_positive_not_cp_ky_fan}

This section proves attainment and the matching upper bound used in
Theorem~\ref{thm:ky_fan_maximum_principle}. For Hermitian $A$, write $S_k(A)$
for the sum introduced in Definition~\ref{def:ky_fan_norm}.

\begin{theorem}[Achievability of the largest-eigenvalue sum]
    \label{thm:ky_fan_achievability}
    \lean{Matrix.IsHermitian.exists_isProj_trace_eq_kyFanNorm}
    \leanok
    \uses{def:ky_fan_norm}
    For every Hermitian $A\in\MN{D}$ and every $k\ge0$, there is an orthogonal
    projection $P$ of rank $\min(k,D)$ with $\re\tr(PA)=S_k(A)$.
\end{theorem}

\begin{proof}\leanok
    \uses{def:ky_fan_norm}
    Diagonalize $A=U\diag(\lambda_i)U^\dagger$ and take $P$ to be the
    projection onto the span of the eigenvectors with the $\min(k,D)$ largest
    eigenvalues. Then $\re\tr(PA)=\sum_{i\le k}\lambda_i=S_k(A)$.
\end{proof}

\begin{theorem}[Upper bound for the largest-eigenvalue sum]
    \label{thm:ky_fan_upper_bound}
    \lean{Matrix.IsHermitian.trace_mul_re_le_kyFanNorm}
    \leanok
    \uses{def:ky_fan_norm}
    For every Hermitian $A\in\MN{D}$ and every orthogonal projection $P$ of
    rank $k<D$, one has $\re\tr(PA)\le S_k(A)$.
\end{theorem}

\begin{proof}\leanok
    \uses{def:ky_fan_norm}
    Conjugating $P$ by $U$ gives an orthogonal projection $Q$ of the same
    rank, and $\re\tr(PA)=\sum_iw_i\lambda_i$, where $w_i$ are the diagonal
    entries of $Q$. Each $w_i$ lies in $[0,1]$ and the weights sum to
    $\tr(P)=k$. With the threshold $c=\lambda_{k+1}$,
    \begin{align}
        \sum_{i\le k}\lambda_i-\sum_iw_i\lambda_i
         & =\sum_{i\le k}(\lambda_i-c)(1-w_i)
        +\sum_{i>k}(c-\lambda_i)w_i
        \ge0,
        \notag
    \end{align}
    since for $i\le k$ the eigenvalues dominate $c$ while $1-w_i\ge0$, and
    for $i>k$ they are dominated by $c$ while $w_i\ge0$.
\end{proof}

Together, Theorems~\ref{thm:ky_fan_achievability} and
\ref{thm:ky_fan_upper_bound} prove Theorem~\ref{thm:ky_fan_maximum_principle},
which supplies the projection estimate in the maximal-overlap theorem.

\section{Right-tensor identities and Choi compressions}
\label{sec:app_positive_not_cp_choi_compression}

This section proves the matrix identities and rank parametrizations used in
the Schmidt-rank Choi criterion, the rectangular compression criterion, and
the rank-$k$ projection criterion. Throughout, $R_X$ and $C_T(X)$ are the
right-tensor factor and Choi compression of
Definitions~\ref{def:choi_right_tensor_factor}
and~\ref{def:choi_right_compression}. Here
$T:\MN{d}\to\MN{d'}$, $X\in M_{d\times k}(\C)$, the Choi matrix acts on
$\C^{d'}\otimes\C^d$, and the compressed matrix acts on
$\C^{d'}\otimes\C^k$.

\begin{theorem}[Right tensor sandwich formula]
    \label{thm:choi_right_tensor_sandwich}
    \lean{ChoiJamiolkowski.rightTensorMatrix_mul_choiMatrix_mul_conjTranspose_rectangular}
    \leanok
    \uses{def:choi_right_compression, def:choi_right_tensor_factor}
    If $\tau$ is the Choi matrix of $T$, then
    \begin{align}
        R_X\tau R_X^\dagger
         & =C_T(X),
        \label{eq:schwarz_right_tensor_sandwich}
    \end{align}
    where $C_T(X)$ is the right-factor Choi compression by $X$.
\end{theorem}

\begin{proof}
    \leanok
    By (\ref{eq:schwarz_right_tensor_factor}), the $(i,p),(j,q)$ entry of the
    left-hand side of (\ref{eq:schwarz_right_tensor_sandwich}) is
    \begin{align}
        \sum_{r,s,a,b}
        \delta_{i,r}X_{a,p}\tau_{(r,a),(s,b)}
        \delta_{j,s}\overline{X_{b,q}}
         & =\sum_{a,b}X_{a,p}\tau_{(i,a),(j,b)}
        \overline{X_{b,q}}.
        \notag
    \end{align}
    which is (\ref{eq:schwarz_right_choi_compression}).
\end{proof}

\begin{lemma}[Compressed maximally entangled vectors have bounded Schmidt rank]
    \label{lem:compressed_omega_schmidt_rank}
    \lean{ChoiJamiolkowski.compressedOmegaVector_hasSchmidtRankLE}
    \leanok
    \uses{def:schmidt_rank, def:choi_right_compression}
    For $X\in M_{d\times k}(\C)$, the vector
    \begin{align}
        \psi_X
         & =\left(d^{-1/2}X_{a,p}\right)_{(a,p)}
        \in\C^d\otimes\C^k
        \label{eq:schwarz_compressed_omega}
    \end{align}
    has Schmidt rank at most $k$.
\end{lemma}

\begin{proof}
    \leanok
    The coefficient matrix of $\psi_X$ is a scalar multiple of $X$. Therefore
    its rank is bounded by the number of columns, which is $k$.
\end{proof}

\begin{theorem}[Rank of a compressed maximally entangled vector]
    \label{thm:compressed_omega_schmidt_rank_eq_rank}
    \lean{ChoiJamiolkowski.compressedOmegaVector_schmidtRank_eq_rank}
    \leanok
    \uses{def:schmidt_rank, def:choi_right_compression}
    Assume $d>0$. For $X\in M_{d\times k}(\C)$, the vector $\psi_X$ in
    (\ref{eq:schwarz_compressed_omega}) satisfies
    $\SR(\psi_X)=\rank X$.
\end{theorem}

\begin{proof}
    \leanok
    The coefficient matrix of $\psi_X$ is $d^{-1/2}X$. Since $d>0$, the
    scalar $d^{-1/2}$ is nonzero, and multiplication by this scalar preserves
    rank.
\end{proof}

\begin{theorem}[Rectangular representative with exact Schmidt rank]
    \label{thm:schmidt_rank_exact_parametrization}
    \lean{ChoiJamiolkowski.exists_compression_of_vector}
    \leanok
    \uses{def:schmidt_rank, def:choi_right_compression,
        thm:compressed_omega_schmidt_rank_eq_rank}
    Assume $d>0$. For every $\psi\in\C^d\otimes\C^k$, there is
    $X\in M_{d\times k}(\C)$ such that
    \begin{align}
        \psi_{(i,p)}
         & =d^{-1/2}X_{i,p}, \notag \\
        \rank X
         & =\SR(\psi).
        \label{eq:schwarz_schmidt_representative}
    \end{align}
    In particular, if $\SR(\psi)\le r$, then $X$ may be chosen with
    $\rank X\le r$.
\end{theorem}

\begin{proof}
    \leanok
    Take $X_{i,p}=d^{1/2}\psi_{(i,p)}$. Since $d>0$, this gives
    $d^{-1/2}X_{i,p}=\psi_{(i,p)}$. The rank identity follows from
    Theorem~\ref{thm:compressed_omega_schmidt_rank_eq_rank}. The bounded-rank
    statement follows by comparison with the assumed Schmidt-rank bound.
\end{proof}

\begin{lemma}[Square parametrization of bounded Schmidt rank]
    \label{lem:schmidt_rank_square_compressed_omega}
    \lean{ChoiJamiolkowski.hasSchmidtRankLE_iff_exists_rank_le_compressedOmegaVector}
    \leanok
    \uses{def:schmidt_rank, def:choi_right_compression,
        thm:schmidt_rank_exact_parametrization}
    Assume $d>0$. A vector $\psi\in\C^d\otimes\C^d$ has Schmidt rank at most
    $k$ if and only if there is a matrix $X\in\MN{d}$ of rank at most $k$
    such that $\psi_{(i,j)}=d^{-1/2}X_{i,j}$.
\end{lemma}

\begin{proof}
    \leanok
    The forward direction is the bounded-rank part of
    Theorem~\ref{thm:schmidt_rank_exact_parametrization}. Conversely,
    if $\psi_{(i,j)}=d^{-1/2}X_{i,j}$, then
    $C_\psi=d^{-1/2}X$, so $\rank(C_\psi)=\rank(X)$ because multiplication
    by a nonzero scalar preserves rank.
\end{proof}

\begin{lemma}[Schmidt rank under the adjoint right tensor factor]
    \label{lem:right_tensor_adjoint_schmidt_rank}
    \lean{ChoiJamiolkowski.schmidtCoeffMatrix_rightTensorMatrix_conjTranspose_mulVec,
        ChoiJamiolkowski.rightTensorMatrix_conjTranspose_mulVec_hasSchmidtRankLE}
    \leanok
    \uses{def:schmidt_rank, def:choi_right_tensor_factor}
    Let $X\in M_{d\times k}(\C)$ and let
    $\eta\in\C^{d'}\otimes\C^k$. Then
    $R_X^\dagger\eta\in\C^{d'}\otimes\C^d$ has Schmidt rank at most $k$.
\end{lemma}

\begin{proof}
    \leanok
    If $\eta$ is read as a $d'\times k$ coefficient matrix $E$, then the
    coefficient matrix of $R_X^\dagger\eta$ is $EX^\dagger$. Therefore
    $\rank(EX^\dagger)\le\rank(X^\dagger)=\rank(X)\le k$.
\end{proof}

\begin{lemma}[Right-tensor representatives of bounded Schmidt rank]
    \label{lem:right_tensor_adjoint_schmidt_rank_representatives}
    \lean{ChoiJamiolkowski.exists_rightTensorMatrix_conjTranspose_mulVec_of_hasSchmidtRankLE}
    \leanok
    \uses{def:schmidt_rank, def:choi_right_tensor_factor,
        lem:matrix_rank_factorization_coordinate}
    If $\psi\in\C^{d'}\otimes\C^d$ has Schmidt rank at most $k$, then there
    exist $X\in M_{d\times k}(\C)$ and
    $\eta\in\C^{d'}\otimes\C^k$ such that
    $\psi=R_X^\dagger\eta$.
\end{lemma}

\begin{proof}
    \leanok
    Let $C_\psi$ be the coefficient matrix of $\psi$. The rank assumption
    gives a factorization $C_\psi=BC$ through $\C^k$. Taking
    $X=C^\dagger$ and reading $B$ as the coefficient matrix of $\eta$ gives
    the required identity, because the coefficient matrix of
    $R_X^\dagger\eta$ is $BX^\dagger=BC$.
\end{proof}

\begin{theorem}[Quadratic form of a right Choi compression]
    \label{thm:choi_right_compression_quadratic_form}
    \lean{ChoiJamiolkowski.rightTensor_choiMatrix_quadraticForm_eq_rectangular,
        ChoiJamiolkowski.rightCompression_quadraticForm_eq_choiMatrix_quadraticForm_rectangular}
    \leanok
    \uses{def:choi_right_compression, def:choi_right_tensor_factor,
        thm:choi_right_tensor_sandwich}
    For $X\in M_{d\times k}(\C)$ and
    $\eta\in\C^{d'}\otimes\C^k$,
    \begin{align}
        \langle\eta,C_T(X)\eta\rangle
         & =\langle R_X^\dagger\eta,\tau_T R_X^\dagger\eta\rangle.
        \label{eq:schwarz_right_compression_quadratic}
    \end{align}
\end{theorem}

\begin{proof}
    \leanok
    Substituting (\ref{eq:schwarz_right_tensor_sandwich}) gives
    \begin{align}
        \langle\eta,R_X\tau_T R_X^\dagger\eta\rangle
         & =\langle R_X^\dagger\eta,\tau_T R_X^\dagger\eta\rangle.
        \notag
    \end{align}
\end{proof}

\begin{lemma}[Rank-one ampliation as Choi compression]
    \label{lem:rank_one_ampliation_choi_compression}
    \lean{ChoiJamiolkowski.nPositiveAmpliation_rankOne_eq_rightCompression}
    \leanok
    \uses{def:blockwise_ampliation, def:choi_right_compression}
    Applying the $k$-fold ampliation of $T$ to the rank-one matrix
    $\ketbra{\psi}{\psi}$, where
    $\psi_{(a,p)}=d^{-1/2}X_{a,p}$, is exactly the right-factor compression
    of the Choi matrix of $T$ by $X$.
\end{lemma}

\begin{proof}
    \leanok
    In entries, both sides are
    \begin{align}
        \sum_{a,b}X_{a,p}
        T\!\left(d^{-1}E_{a,b}\right)_{i,j}
        \overline{X_{b,q}}.
        \notag
    \end{align}
    If $\psi_{(a,p)}=d^{-1/2}X_{a,p}$, then for fixed $p,q$, the corresponding
    block of $\ketbra{\psi}{\psi}$ is
    \begin{align}
        (\ketbra{\psi}{\psi})_{(\cdot,p),(\cdot,q)}
         & =d^{-1}\sum_{a,b}X_{a,p}\overline{X_{b,q}}E_{a,b}.
        \notag
    \end{align}
    Hence
    \begin{align}
        T\!\left(d^{-1}\sum_{a,b}
           X_{a,p}\overline{X_{b,q}}E_{a,b}\right)_{i,j}
         & =d^{-1}\sum_{a,b}X_{a,p}\overline{X_{b,q}}
                                   (T(E_{a,b}))_{i,j} \notag \\
         & =\sum_{a,b}X_{a,p}\tau_{(i,a),(j,b)}
        \overline{X_{b,q}}.
        \notag
    \end{align}
    where the second equality uses
    $\tau_{(i,a),(j,b)}=d^{-1}(T(E_{a,b}))_{i,j}$.
\end{proof}

\begin{lemma}[Right-tensor multiplication]
    \label{lem:right_tensor_multiplication}
    \lean{ChoiJamiolkowski.rightTensorMatrix_mul_rightTensorMatrix}
    \leanok
    \uses{def:choi_right_tensor_factor}
    Let $X\in M_{d\times k}(\C)$ and $P\in\MN{d}$. Then the right-tensor
    factors satisfy $R_XR_P=R_{PX}$.
\end{lemma}

\begin{proof}
    \leanok
    This is an entrywise calculation: the Kronecker delta in the physical
    index forces the two right-tensor factors to have the same physical
    coordinate, and the remaining sum is the matrix product $PX$.
\end{proof}

\begin{lemma}[Square sandwich implies fixed rectangular compression]
    \label{lem:square_sandwich_fixed_rectangular_compression}
    \lean{ChoiJamiolkowski.rightCompression_posSemidef_of_projection_sandwich_of_mul_eq_self_rectangular}
    \leanok
    \uses{thm:choi_right_tensor_sandwich, lem:right_tensor_multiplication}
    Let $T:\MN{d}\to\MN{d'}$ have Choi matrix $\tau$. Suppose
    $P\in\MN{d}$ and $X\in M_{d\times k}(\C)$ satisfy $PX=X$. If
    $R_P\tau R_P^\dagger\ge0$, then the rectangular right compression of
    $\tau$ by $X$ is positive semidefinite.
\end{lemma}

\begin{proof}
    \leanok
    Since $PX=X$, Lemma~\ref{lem:right_tensor_multiplication} gives
    $R_XR_P=R_X$. Therefore
    \begin{align}
        R_X\tau R_X^\dagger
         & =R_X(R_P\tau R_P^\dagger)R_X^\dagger,
        \label{eq:schwarz_fixed_rectangular_compression}
    \end{align}
    which is positive because positive semidefiniteness is preserved under
    compression. The left-hand side of
    (\ref{eq:schwarz_fixed_rectangular_compression}) is the rectangular
    compression by (\ref{eq:schwarz_right_tensor_sandwich}).
\end{proof}

\begin{lemma}[Rank-$k$ projection fixing a rectangular right factor]
    \label{lem:rank_k_projection_fixing_rectangular_factor}
    \lean{Matrix.exists_isHermitian_idempotent_rank_mul_eq_self}
    \leanok
    Let $X\in M_{d\times k}(\C)$ and suppose $k\le d$. Then there exists a
    Hermitian projection $P\in\MN{d}$ of rank $k$ such that $PX=X$.
\end{lemma}

\begin{proof}
    \leanok
    The column space of $X$ has dimension at most $k$. Since $k\le d$, extend
    this column space to a $k$-dimensional subspace $U\subseteq\C^d$. Let $P$
    be the orthogonal projection onto $U$. Then $P$ is Hermitian and
    idempotent, its rank is $\dim U=k$, and it fixes each column of $X$.
\end{proof}

\begin{lemma}[Right-tensor factorization for $VV^\dagger$]
    \label{lem:right_tensor_factorization_mul_adjoint}
    \lean{ChoiJamiolkowski.rightTensorMatrix_mul_conjTranspose_eq_conjTranspose_mul_self}
    \leanok
    \uses{def:choi_right_tensor_factor}
    For every $V\in M_{d\times k}(\C)$, the right-tensor factor satisfies
    $R_{VV^\dagger}=R_V^\dagger R_V$.
\end{lemma}

\begin{proof}
    \leanok
    This is the entrywise calculation obtained by expanding the definition of
    $R_V$ and summing over the middle physical index.
\end{proof}

\begin{lemma}[Factorized right-tensor Choi sandwich]
    \label{lem:factorized_right_tensor_choi_sandwich}
    \lean{ChoiJamiolkowski.rightTensorMatrix_mul_conjTranspose_choi_sandwich_eq_rectangular}
    \leanok
    \uses{def:choi_right_tensor_factor,
        lem:right_tensor_factorization_mul_adjoint}
    Let $T:\MN{d}\to\MN{d'}$ have Choi matrix $\tau$. For every
    $V\in M_{d\times k}(\C)$,
    \begin{align}
        R_{VV^\dagger}\tau R_{VV^\dagger}^\dagger
         & =R_V^\dagger(R_V\tau R_V^\dagger)R_V.
        \label{eq:schwarz_factorized_choi_sandwich}
    \end{align}
\end{lemma}

\begin{proof}
    \leanok
    Replace $R_{VV^\dagger}$ by $R_V^\dagger R_V$ using
    Lemma~\ref{lem:right_tensor_factorization_mul_adjoint} and reassociate the
    matrix products.
\end{proof}

\begin{theorem}[Factorized right projections in the Choi criterion]
    \label{thm:factorized_right_projection_choi_sandwich}
    \lean{ChoiJamiolkowski.IsNPositiveMap.rightTensor_choiMatrix_sandwich_posSemidef_of_mul_conjTranspose_rectangular}
    \leanok
    \uses{lem:factorized_right_tensor_choi_sandwich,
        thm:k_positive_right_tensor_choi_sandwiches}
    Assume $d>0$. Let $T:\MN{d}\to\MN{d'}$ be $k$-positive,
    with Choi matrix $\tau$. If a right-factor matrix has the form
    $P=VV^\dagger$ with $V\in M_{d\times k}(\C)$, then
    $R_P\tau R_P^\dagger\geq0$.
\end{theorem}

\begin{proof}
    \leanok
    By (\ref{eq:schwarz_factorized_choi_sandwich}),
    $R_{VV^\dagger}\tau R_{VV^\dagger}^\dagger
        =R_V^\dagger(R_V\tau R_V^\dagger)R_V$.
    The factor $R_V\tau R_V^\dagger$ is positive by
    Theorem~\ref{thm:k_positive_right_tensor_choi_sandwiches}, and positive
    semidefiniteness is preserved under compression by $R_V$.
\end{proof}

\begin{lemma}[Rank-$k$ Hermitian projection factorization]
    \label{lem:rank_k_hermitian_projection_factorization}
    \lean{Matrix.exists_mul_conjTranspose_of_isHermitian_idempotent_rank}
    \leanok
    Let $P\in\MN{D,D}(\C)$ be Hermitian, idempotent, and of rank $k$.
    Then there exists $V\in\MN{D,k}(\C)$ such that $P=VV^\dagger$.
\end{lemma}

\begin{proof}
    \leanok
    Regard $P$ as a linear operator on $\C^D$. Hermiticity and idempotence say
    that this operator is the orthogonal projection onto its range. Choose an
    orthonormal basis of the range, indexed by $\{0,\ldots,k-1\}$ using the
    rank hypothesis, and let $V$ be the matrix whose columns are these basis
    vectors. The standard rank-one expansion of an orthogonal projection gives
    $P=VV^\dagger$.
\end{proof}

These identities provide the bounded-Schmidt-rank representatives used in
Theorems~\ref{thm:k_positive_schmidt_rank_choi_expectation}
and~\ref{thm:schmidt_rank_choi_test_converse}, identify ampliations with Choi
compressions in Theorem~\ref{thm:k_positive_iff_right_choi_compressions}, and
supply the projection factorizations in
Theorem~\ref{thm:rank_k_right_projection_choi_criterion}.

\section{Closure properties of $k$-positive maps}
\label{sec:app_positive_not_cp_k_positive_closure}

The next results record the monotonicity and elementary convex-cone properties
used with the positivity hierarchy in Chapter~\ref{ch:positive_not_cp}.

\begin{theorem}[$(k\!+\!1)$-positive implies $k$-positive]
    \label{thm:k_positive_mono_succ}
    \lean{IsNPositiveMap.mono}
    \leanok
    Positivity under amplification is monotone in $k$.
\end{theorem}

\begin{proof}
    \leanok
    Embed $\MN{k}$ as the upper-left corner of $\MN{k+1}$. For
    $X\succeq0$ in $\MN{D}\otimes\MN{k}$, extend $X$ by a zero row and
    column, apply the positive $(k+1)$-fold ampliation, and compress back to
    the same corner. Compression preserves positive semidefiniteness, and the
    compressed image is the $k$-fold ampliation of $X$.
\end{proof}

\begin{theorem}[$k$-positive implies $m$-positive for $m \le k$]
    \label{thm:k_positive_mono}
    \lean{IsNPositiveMap.mono_of_le}
    \leanok
    If a linear map is $k$-positive and $m \le k$, then it is $m$-positive.
\end{theorem}

\begin{proof}
    \leanok
    Repeat the corner-extension and compression argument from dimension $k$
    to dimension $m$. Equivalently, iterate the preceding theorem $k-m$
    times.
\end{proof}

\begin{theorem}[Zero map is $k$-positive]
    \label{thm:k_positive_zero}
    \lean{IsNPositiveMap.zero}
    \leanok
    The zero map is $k$-positive.
\end{theorem}

\begin{proof}
    \leanok
    Every ampliation of the zero map is zero, and the zero matrix is positive
    semidefinite.
\end{proof}

\begin{theorem}[Sums of $k$-positive maps]
    \label{thm:k_positive_add}
    \lean{IsNPositiveMap.add}
    \leanok
    If $E$ and $F$ are $k$-positive, then $E+F$ is $k$-positive.
\end{theorem}

\begin{proof}
    \leanok
    For $X\succeq0$, both $E^{(k)}(X)$ and $F^{(k)}(X)$ are positive
    semidefinite. Since $(E+F)^{(k)}=E^{(k)}+F^{(k)}$, their sum is positive
    semidefinite.
\end{proof}

\begin{theorem}[Nonnegative multiples of $k$-positive maps]
    \label{thm:k_positive_smul_nonneg}
    \lean{IsNPositiveMap.smul_nonneg}
    \leanok
    If $E$ is $k$-positive and $c \ge 0$, then $cE$ is $k$-positive.
\end{theorem}

\begin{proof}
    \leanok
    Ampliation commutes with scalar multiplication, and a non-negative scalar
    multiple of a positive semidefinite matrix is positive semidefinite.
\end{proof}

\begin{theorem}[Closedness of $k$-positive maps]
    \label{thm:k_positive_closed}
    \lean{isClosed_setOf_isNPositiveMap}
    \leanok
    For each $k$, the set of $k$-positive linear maps is closed in the
    finite-dimensional space of linear maps.
\end{theorem}

\begin{proof}
    \leanok
    Let $E_j\to E$ with every $E_j$ $k$-positive. For each fixed
    $X\succeq0$, continuity of ampliation and evaluation gives
    $E_j^{(k)}(X)\to E^{(k)}(X)$. The positive semidefinite cone is closed, so
    $E^{(k)}(X)\succeq0$. Since this holds for every positive semidefinite
    $X$, the limit map $E$ is $k$-positive.
\end{proof}

\section{Schmidt-number geometry and entanglement witnesses}
\label{sec:app_positive_not_cp_schmidt_number_geometry}

This section supplies the convexity, compactness, and separation results used
in the witness criterion of Theorem~\ref{thm:not_has_schmidt_number_le_iff_exists_witness}.

\begin{theorem}[Basic properties of bounded Schmidt number]
    \label{thm:has_schmidt_number_le_basic}
    \lean{Matrix.HasSchmidtNumberLE.posSemidef, Matrix.HasSchmidtNumberLE.add,
        Matrix.HasSchmidtNumberLE.mono, Matrix.HasSchmidtNumberLE.smul}
    \leanok
    \uses{def:has_schmidt_number_le}
    A matrix of Schmidt number at most $n$ is positive semidefinite. States of
    Schmidt number at most $n$ are closed under addition and under multiplication by a
    non-negative real scalar, and a bound on the Schmidt number relaxes to any larger
    bound.
\end{theorem}

\begin{proof}
    \leanok
    Each summand $\ket{\psi_i}\bra{\psi_i}$ is a rank-one positive semidefinite
    matrix, and a finite sum of positive semidefinite matrices is positive
    semidefinite. Concatenating two pure-state families realizes a sum as a single
    finite sum of pure-state projectors of Schmidt rank at most $n$, and a pure
    summand of Schmidt rank at most $n$ also has Schmidt rank at most any larger
    bound. Multiplication by a non-negative real $a$ rescales each pure summand
    $\psi_i$ to $\sqrt a\,\psi_i$, which leaves its Schmidt rank unchanged and
    reproduces $a$ times the original state.
\end{proof}

\begin{theorem}[Convexity of the Schmidt-number set]
    \label{thm:convex_setOf_has_schmidt_number_le}
    \lean{Matrix.convex_setOf_hasSchmidtNumberLE}
    \leanok
    \uses{def:has_schmidt_number_le}
    For each $n$ the set $S_n$ of bipartite states of Schmidt number at most $n$ is
    convex.
\end{theorem}

\begin{proof}
    \leanok
    A convex combination $a x + b y$ with $a,b\ge 0$ is a sum of two non-negatively
    scaled states of Schmidt number at most $n$, and that bound is preserved under
    non-negative scaling and under addition; rescaling a vector by a non-negative
    scalar scales its coefficient matrix and so does not increase the Schmidt rank.
\end{proof}

\begin{theorem}[Compactness of the Schmidt-number state set]
    \label{thm:compact_setOf_has_schmidt_number_le_trace_one}
    \lean{Matrix.isCompact_setOf_hasSchmidtNumberLE_trace_one}
    \leanok
    \uses{def:has_schmidt_number_le, def:schmidt_rank,
        thm:convex_setOf_has_schmidt_number_le}
    For each $n$ the set $S_n$ of trace-one bipartite states of Schmidt number at most
    $n$ is compact.
\end{theorem}

\begin{proof}
    \leanok
    A trace-one state of Schmidt number at most $n$ is a convex combination of pure
    states $\ket{\psi}\bra{\psi}$ whose vector $\psi$ has unit norm and Schmidt rank
    at most $n$, so $S_n$ is the convex hull of the set $P_n$ of those pure states.
    The set $P_n$ is the image, under the continuous projector map
    $\psi\mapsto\ket{\psi}\bra{\psi}$, of the unit vectors of Schmidt rank at most
    $n$. In the Euclidean norm those unit vectors form the unit sphere, which is
    compact in finite dimension, intersected with the closed set on which the Schmidt
    rank is at most $n$ (the coefficient matrix depends continuously on the vector and
    the matrices of rank at most $n$ form a closed set); hence $P_n$ is compact. The
    trace of $\ket{\psi}\bra{\psi}$ is the squared Euclidean norm of $\psi$, so the
    trace-one constraint is exactly the unit-norm constraint and the weights of the
    convex decomposition are the squared norms, summing to the trace. The convex hull
    of a compact set is compact in a finite-dimensional real normed space.
\end{proof}

\begin{theorem}[Trace-form representation of a real-linear functional]
    \label{thm:exists_isHermitian_trace_form_re}
    \lean{Matrix.exists_isHermitian_trace_form_re}
    \leanok
    Every real-linear functional $g$ on the space of square complex matrices is the
    trace form of a Hermitian matrix: there is a Hermitian $H_0$ with
    $g(X)=\re\tr(X H_0)$ for every Hermitian $X$.
\end{theorem}

\begin{proof}
    \leanok
    The real bilinear pairing $(X,H)\mapsto\re\tr(X H)$ is nondegenerate,
    because on the diagonal pair $(H,H^{\dagger})$ it equals the squared Frobenius norm
    $\re\tr(H H^{\dagger})=\sum_{i,j}\lvert H_{ij}\rvert^2$, which is positive
    unless $H=0$. Hence the map sending $H$ to the functional
    $X\mapsto\re\tr(X H)$ is an injective real-linear map from the matrix
    space to its real dual. These two spaces have the same finite real dimension, so the
    map is also surjective, and every $g$ is represented by some $H$. Replacing $H$ by
    its Hermitian part $\tfrac12(H+H^{\dagger})$ keeps the value on Hermitian inputs,
    since $\tr(X H^{\dagger})=\overline{\tr(X H)}$ when $X$ is Hermitian, and produces a
    Hermitian representing matrix.
\end{proof}

\begin{theorem}[Entanglement witness from a separating hyperplane]
    \label{thm:exists_entanglement_witness}
    \lean{Matrix.exists_isHermitian_witness}
    \leanok
    \uses{def:has_schmidt_number_le, def:schmidt_rank,
        thm:convex_setOf_has_schmidt_number_le,
        thm:compact_setOf_has_schmidt_number_le_trace_one,
        thm:exists_isHermitian_trace_form_re}
    Let $\rho$ be a trace-one Hermitian bipartite state whose Schmidt number exceeds $n$.
    Then there is a Hermitian operator $W$ such that, for every $\psi$ of Schmidt rank
    at most $n$,
    \begin{align}
        \re\tr(W\rho)            & <0,
        \notag                            \\
        \re\bra{\psi}W\ket{\psi} & \ge 0.
        \notag
    \end{align}
    This is the only-if direction of Wolf's Proposition~3.3
    \cite[Chapter~3, Proposition~3.3]{Wolf2012Quantum}: a state of Schmidt number larger
    than $n$ is detected by an entanglement witness for $S_n$.
\end{theorem}

\begin{proof}
    \leanok
    The set $S_n$ of trace-one states of Schmidt number at most $n$ is convex and
    compact, and $\rho\notin S_n$. The geometric Hahn--Banach theorem separates the
    closed convex set $\{\rho\}$ from the compact convex set $S_n$ by a continuous
    real-linear functional $f$ and a constant $c$ with $f(\rho)<c<f(\sigma)$ for every
    $\sigma\in S_n$. Representing $f$ as the trace form of a Hermitian $H_0$ and setting
    $W=H_0-c\,\Id$ gives $\re\tr(W\rho)=f(\rho)-c<0$. For a vector $\psi$
    of Schmidt rank at most $n$ the projector $\ket{\psi}\bra{\psi}$, after
    normalization by $\lVert\psi\rVert^2$, is a trace-one state of $S_n$, so its value
    under $f$ is at least $c$; multiplying back by the squared norm gives
    $\re\bra{\psi}W\ket{\psi}\ge 0$, with the zero vector giving the value
    $0$ directly.
\end{proof}

\section{Entanglement witnesses and Choi trace pairing}
\label{sec:app_positive_not_cp_witness_choi}

The following results identify witnesses with Choi matrices of $n$-positive
maps and supply the trace identity used in
Theorem~\ref{thm:exists_isNPositiveMap_detects}.

\begin{theorem}[Surjectivity of the Choi correspondence]
    \label{thm:exists_choiMatrix_eq}
    \lean{ChoiJamiolkowski.exists_choiMatrix_eq}
    \leanok
    \uses{thm:choi_rect_mutual_inverse}
    Assume $D>0$. Every bipartite matrix $W$ on $\C^D\otimes\C^D$ is the Choi matrix
    $\tau_T=(T\otimes\Id)(\ket{\Omega}\bra{\Omega})$ of some linear map
    $T\colon\MN{D}\to\MN{D}$.
\end{theorem}

\begin{proof}
    \leanok
    Apply the explicit inverse $\tau\mapsto T$ from the rectangular Choi correspondence
    (Theorem~\ref{thm:choi_rect_mutual_inverse}) to $W$. Its inverse law gives
    $\tau_T=W$; the displayed square statement is the specialization $d=d'=D$.
\end{proof}

\begin{theorem}[A witness is the Choi matrix of an $n$-positive map]
    \label{thm:exists_isNPositiveMap_choiMatrix_eq}
    \lean{ChoiJamiolkowski.exists_isNPositiveMap_choiMatrix_eq_of_witness}
    \leanok
    \uses{def:k_positive_map, def:schmidt_rank, thm:choi_rect_mutual_inverse,
        thm:schmidt_rank_choi_test_iff, thm:exists_entanglement_witness}
    Assume $d'>0$. Let $W$ be a Hermitian operator on $\C^d\otimes\C^{d'}$ with
    $\re\bra{\psi}W\ket{\psi}\ge 0$ for every $\psi$ of Schmidt rank at most $n$.
    Then $W$ is the Choi matrix of an $n$-positive map
    $P\colon\MN{d'}\to\MN{d}$. This is the
    Choi--Jamio\l kowski translation of the entanglement witness of Wolf's Proposition~3.3
    into the $n$-positive map of Wolf's Proposition~3.4
    \cite[Chapter~3, Proposition~3.4]{Wolf2012Quantum}.
\end{theorem}

\begin{proof}
    \leanok
    The explicit rectangular inverse gives $P\colon\MN{d'}\to\MN{d}$ with $\tau_P=W$.
    Because $W$ is Hermitian, the Choi quadratic form
    $\bra{\psi}W\ket{\psi}$ is real, and it equals the witness expectation
    $\tr(W\ket{\psi}\bra{\psi})$; the witness condition makes it non-negative on every
    vector of Schmidt rank at most $n$. This is exactly the Schmidt-rank Choi criterion
    (\ref{thm:schmidt_rank_choi_test_iff}) for $n$-positivity of $P$.
\end{proof}

\begin{theorem}[Choi trace pairing through the trace-pairing adjoint]
    \label{thm:trace_choiMatrix_omega_quadratic_adjoint}
    \lean{ChoiJamiolkowski.trace_choiMatrix_mul_eq_omegaVec_quadraticForm_traceAdjointMap}
    \leanok
    \uses{def:trace_pairing_adjoint, thm:trace_pairing_adjoint_ampliation}
    For a linear map $P\colon\MN{d'}\to\MN{d}$ with Choi matrix $\tau_P$ and any bipartite
    matrix $\rho$ on $\C^d\otimes\C^{d'}$,
    \begin{align}
        \tr(\tau_P\,\rho)
         & =\bra{\Omega_{d'}}(P^{*}\otimes\Id_{d'})(\rho)\ket{\Omega_{d'}},
        \notag
    \end{align}
    where $P^{*}\colon\MN{d}\to\MN{d'}$ is the trace-pairing adjoint of $P$.
\end{theorem}

\begin{proof}
    \leanok
    Since $\tau_P=(P\otimes\Id_{d'})(\ket{\Omega_{d'}}\bra{\Omega_{d'}})$, cyclicity
    of the trace gives
    $\tr(\tau_P\rho)=\tr(\rho\,(P\otimes\Id_{d'})(\ket{\Omega_{d'}}\bra{\Omega_{d'}}))$.
    The
    $\Id$-ampliation of the trace-pairing adjoint is the trace-pairing adjoint of the
    ampliation, so moving $P$ across the trace pairing replaces
    $(P\otimes\Id_{d'})$ by $(P^{*}\otimes\Id_{d'})$ acting on $\rho$ and turns the
    pairing against $\ket{\Omega_{d'}}\bra{\Omega_{d'}}$ into the quadratic form
    $\bra{\Omega_{d'}} \cdot\,\ket{\Omega_{d'}}$.
\end{proof}
